\VerdierPhysicalPageAnchor{0039}%
\VerdierSemanticAnchor{ch1-1-2-8}
\noindent\textbf{Definition 1.2.8.}\quad
Let $(\mathcal C_e,G^\circ,T,c)$ and
$(\mathcal C'_e,G^\circ,T',c')$ be two $G^\circ$-categories, and let
$(F_e,m)$ and $(F'_e,m')\colon\mathcal C_e\to\mathcal C'_e$ be two
$G^\circ$-functors. A morphism of functors
$u_e\colon F_e\to F'_e$ is said to be \emph{compatible with the actions
of $G^\circ$} if, for every $g\in G$, diagram (1.2.7.2) commutes.

\VerdierSemanticAnchor{ch1-1-2-additive}
Finally, we shall say that a $G^\circ$-category is \emph{additive} if
its underlying category is additive. By (1.1.3), the
$G^\circ$-category associated with an additive graded category is an
additive $G^\circ$-category.

\VerdierSemanticAnchor{ch1-1-3}
\subsection{Equivalence of Two Languages}

\VerdierSemanticAnchor{ch1-1-3-1}
\noindent\textbf{1.3.1.}\quad
Fix a universe $U$. Let $\operatorname{Grad}(G)$ denote the
$U$-category whose objects are the $G$-graded categories belonging to
$U$ and whose morphisms are graded functors. Likewise, let
$G^\circ\text{-}\operatorname{Cat}$ denote the $U$-category whose
objects are the $G^\circ$-categories belonging to $U$ and whose
morphisms are $G^\circ$-functors. It follows immediately from (1.2.1),
(1.2.3), and (1.2.5) that choosing a normalized cleavage for every
$G$-graded category defines a functor

\VerdierSemanticAnchor{ch1-1-3-1-eq1}
\begin{equation}
  \Psi\colon \operatorname{Grad}(G)
  \longrightarrow G^\circ\text{-}\operatorname{Cat}.
  \tag{1.3.1.1}
\end{equation}

The result of [18] can be interpreted as follows.

\VerdierSemanticAnchor{ch1-1-3-2}
\noindent\textbf{Proposition 1.3.2.}\quad
\emph{The functor $\Psi$ is an equivalence of categories.}

\VerdierSemanticAnchor{ch1-1-3-2-quasi-inverse}
A description of a quasi-inverse functor to $\Psi$ may be found in
[18]. Here we shall only give a rapid description of the graded
category $\mathcal C[G]$ associated with a $G^\circ$-category
$(\mathcal C,G^\circ,T,c)$. First, the set of objects of
$\mathcal C[G]$ is the set of objects of $\mathcal C$. Let $X$ and $Y$
be two objects of $\mathcal C$. Set

\VerdierSemanticAnchor{ch1-1-3-2-eq1}
\begin{equation}
  \operatorname{Hom}^{g}(X,Y)
  =\operatorname{Hom}_{\mathcal C}(X,T(g)Y),
  \qquad g\in G,
  \tag{1.3.2.1}
\end{equation}

\VerdierSemanticAnchor{ch1-1-3-2-eq2}
\begin{equation}
  \operatorname{Hom}_{\mathcal C[G]}(X,Y)
  =\coprod_{g\in G}\operatorname{Hom}^{g}(X,Y).
  \tag{1.3.2.2}
\end{equation}

\VerdierSemanticAnchor{ch1-1-3-2-composition}
We now describe the composition of morphisms in $\mathcal C[G]$.
Let $X$, $Y$, and $Z$ be three objects of $\mathcal C$, and let
$u\in\operatorname{Hom}^{g}(X,Y)$ and
$v\in\operatorname{Hom}^{g'}(Y,Z)$, where $g,g'\in G$. The functor
$T(g)$ defines a map

\VerdierSemanticAnchor{ch1-1-3-2-eq3}
\begin{equation}
  \operatorname{Hom}^{g'}(Y,Z)
  \longrightarrow
  \operatorname{Hom}_{\mathcal C}
    \bigl(T(g)Y,T(g)T(g')Z\bigr)
  \tag{1.3.2.3}
\end{equation}
